\documentclass[12pt]{article}
\usepackage[utf8]{inputenc}
\usepackage[T1]{fontenc}
\usepackage[margin=1in]{geometry}
\usepackage{setspace}
\usepackage{xcolor}
\usepackage[colorlinks=true, urlcolor=blue!60!black, linkcolor=blue!60!black]{hyperref}
\setlength{\parindent}{0em}
\setlength{\parskip}{0.45em}
\renewcommand{\baselinestretch}{1.0}
\begin{document}
\begin{center}
{\Large \textit{Cuyler v.\ Adams}, 449 U.S.\ 433 (1981)\par}
\vspace{0.3em}
{\normalsize Case Brief --- One-Page Edition\par}
\end{center}

\textbf{Memory Jogger.} A congressionally consented compact is federal law; the Interstate Agreement on Detainers preserves the right to an extradition hearing (\href{https://supreme.justia.com/cases/federal/us/449/433/}{Justia}).

\textbf{Facts.} John Adams, imprisoned in Pennsylvania, was transferred to New Jersey for trial under Article IV of the Interstate Agreement on Detainers without the pre-transfer hearing the Uniform Criminal Extradition Act guarantees.

\textbf{Procedural History.} His \S~1983 suit was dismissed; the Third Circuit held the IAD federal law and reversed.

\textbf{Votes.} Affirmed 6--3: Brennan for the Court; Rehnquist, with Burger and Stewart, dissented.

\textbf{Holding.} Congressional consent, even in advance, to a compact on a subject appropriate for federal legislation makes it federal law. Article IV(d) preserves the prisoner's pre-transfer hearing.

\textbf{Reasoning.} Consent is required only for agreements threatening federal supremacy, but consent actually given federalizes the compact; the 1934 Crime Control Consent Act covered the IAD. Article IV(d)'s ``any right'' language and single exception preserve other safeguards. Rehnquist called it ``judicial alchemy'' and warned of a ``judicial Midas.''

\textbf{Examples.} \emph{Same-Side:} a consented disaster-response compact. \emph{Opposite-Side:} an unconsented licensing compact. \emph{Fence-Sitter:} the Uniform Act to Secure the Attendance of Witnesses.

\textbf{Critique.} Treating a blanket advance consent as federal lawmaking blurs permission and enactment, and the Court left unclear which reciprocal state statutes count as compacts.

\textbf{Key Quotation.} ``But where Congress has authorized the States to enter into a cooperative agreement, and where the subject matter of that agreement is an appropriate subject for congressional legislation, the consent of Congress transforms the States' agreement into federal law under the Compact Clause.'' --- Justice Brennan, 449 U.S.\ at 440.

\end{document}
